\documentclass[aspectratio=169]{beamer}
\usetheme{Madrid}
\usecolortheme{default}
\usepackage[utf8]{inputenc}
\usepackage[T5]{fontenc}
\usepackage{graphicx}
\usepackage{hyperref}
\usepackage{booktabs}

\setbeamertemplate{caption}[numbered]
\renewcommand{\figurename}{Hình}

\title[Trí tuệ Nhân tạo cho Kế toán]{Trí tuệ Nhân tạo cho Kế toán \\ \vspace{0.3cm} \Large KẾ HOẠCH SLIDE LÝ THUYẾT - DAY 07 (BUỔI 7)}
\author{Đại học Đông Á}
\date{\today}

\begin{document}

% SLIDE 1
\begin{frame}
    \titlepage
\end{frame}

% SLIDE 2
\begin{frame}{Nội dung Bài học}
    \tableofcontents
\end{frame}


\section{KẾ HOẠCH SLIDE LÝ THUYẾT - DAY 07 (BUỔI 7)}


\begin{frame}{Title: Buổi 7: Ứng dụng AI hỗ trợ lập báo cáo theo yêu cầu (Dashboards).}
    \begin{itemize}
        \item Phụ đề: Chương 9 - Communicating Data Analysis Results
        \item Giảng viên: Đại học Đông Á
        \item Môn học: Trí tuệ Nhân tạo cho Kế toán
    \end{itemize}
\end{frame}


\begin{frame}{Năng lực đạt được sau buổi học.}
    \begin{itemize}
        \item \textbf{Về Kiến thức:} Hiểu được vai trò của Data Storytelling trong kế toán và các bước tạo ra Data Visualization (Dashboards) hiệu quả.
        \item \textbf{Về Kỹ năng:} Nhận diện được các biểu đồ gây hiểu lầm (Misleading Data Visualizations) và biết cách lập báo cáo, thuyết minh BCTC có sự hỗ trợ của AI.
        \item \textbf{Về Tư duy:} Đánh giá phê phán các đề xuất của AI khi trực quan hóa, lựa chọn đúng loại biểu đồ cho từng mục tiêu giao tiếp.
    \end{itemize}
\end{frame}


\begin{frame}{Nội dung Chương trình (Agenda).}
    \begin{itemize}
        \item 1. Kể chuyện bằng dữ liệu (Data Storytelling).
        \item 2. Các bước tạo Trực quan hóa dữ liệu hiệu quả.
        \item 3. Đặc điểm của một Data Visualization tốt.
        \item 4. Nhận diện các biểu đồ sai lệch (Misleading Visualizations).
        \item 5. Dashboard và Trực quan hóa tương tác.
        ### Phần 1: Kể chuyện bằng dữ liệu (Data Storytelling)
    \end{itemize}
\end{frame}


\begin{frame}{Data Literacy - Hiểu biết về Dữ liệu.}
    \begin{itemize}
        \item Khái niệm: Là khả năng đọc, hiểu, sáng tạo và giao tiếp bằng dữ liệu.
        \item Tầm quan trọng trong Kế toán: Không chỉ báo cáo số liệu mà phải hiểu được ý nghĩa đằng sau những con số.
    \end{itemize}
\end{frame}


\begin{frame}{Giao tiếp hiệu quả trong Kế toán (Communicate Effectively).}
    \begin{itemize}
        \item Hiểu đối tượng người nghe (Audience).
        \item Tập trung vào thông điệp cốt lõi (Message).
        \item Đặt dữ liệu vào đúng bối cảnh (Context).
    \end{itemize}
\end{frame}


\begin{frame}{Đối tượng người nghe (The Audience).}
    \begin{itemize}
        \item Stakeholders nội bộ (Internal): Đã có kiến thức nền về công ty.
        \item Stakeholders bên ngoài (External - Nhà đầu tư, ngân hàng): Cần thông tin tóm lược, rõ ràng và có bối cảnh.
    \end{itemize}
\end{frame}


\begin{frame}{Yếu tố cấu thành một Data Story.}
    \begin{itemize}
        \item Gồm 3 yếu tố: Dữ liệu (Data), Tường thuật (Narrative), và Hình ảnh (Visuals).
        \item Data + Narrative = Explain (Giải thích).
        \item Data + Visuals = Enlighten (Làm sáng tỏ).
        \item Narrative + Visuals = Engage (Thu hút).
    \end{itemize}
\end{frame}


\begin{frame}{Cấu trúc của một Data Story (Freytag's Pyramid).}
    \begin{itemize}
        \item Mở bài (Exposition): Giới thiệu vấn đề.
        \item Thắt nút (Rising Action): Khám phá chi tiết dữ liệu.
        \item Cao trào (Climax): Điểm mấu chốt, Aha moment.
        \item Mở nút (Falling Action) \& Kết luận (Resolution).
    \end{itemize}
\end{frame}


\begin{frame}{Ví dụ về Data Storytelling trong Kế toán.}
    \begin{itemize}
        \item Khởi đầu: Nghi ngờ gian lận chi phí mua hàng.
        \item Cao trào: Trực quan hóa chỉ ra sự bất thường ở một nhà cung cấp cụ thể.
        \item Kết luận: Đề xuất siết chặt kiểm soát nội bộ.
        ### Phần 2: Các bước tạo Trực quan hóa dữ liệu hiệu quả
    \end{itemize}
\end{frame}


\begin{frame}{Data Visualization là gì?}
    \begin{itemize}
        \item Là quá trình hiển thị dữ liệu để mang lại ý nghĩa và insights cho người xem.
        \item Giúp chuyển đổi bảng số liệu khô khan (BCTC) thành hình ảnh sinh động.
    \end{itemize}
\end{frame}


\begin{frame}{Bước 1: Xác minh Dữ liệu (Verify the Data).}
    \begin{itemize}
        \item Tránh nguyên tắc Garbage in, garbage out.
        \item Dữ liệu phải Chính xác (Accurate) và Đầy đủ (Complete).
    \end{itemize}
\end{frame}


\begin{frame}{Tính Nhất quán và Cập nhật của Dữ liệu.}
    \begin{itemize}
        \item Nhất quán (Consistent): Cùng định dạng, cùng cấp độ chi tiết qua các kỳ.
        \item Cập nhật (Fresh \& Timely): Dữ liệu phải là mới nhất cho kỳ báo cáo.
    \end{itemize}
\end{frame}


\begin{frame}{Bước 2: Xem xét Đối tượng khán giả (Consider the Audience).}
    \begin{itemize}
        \item Novice (Người mới): Cần nhiều thông tin nền tảng, giải thích chi tiết.
        \item Managerial (Quản lý): Cần kết quả có thể hành động ngay (Actionable results).
    \end{itemize}
\end{frame}


\begin{frame}{Đối tượng Chuyên gia và Điều hành.}
    \begin{itemize}
        \item Expert (Chuyên gia): Muốn khám phá sâu, công cụ tương tác mạnh (Drill-down).
        \item Executive (Điều hành): Cần thông tin vĩ mô, những insights quan trọng nhất.
    \end{itemize}
\end{frame}


\begin{frame}{Bước 3: Xác định Mục tiêu (Define the Objective).}
    \begin{itemize}
        \item Cần trả lời câu hỏi/vấn đề gì?
        \item Mục tiêu có thể là: Thể hiện thành phần (Composition), Mối quan hệ (Relationships)...
    \end{itemize}
\end{frame}


\begin{frame}{Các loại Mục tiêu trực quan hóa khác.}
    \begin{itemize}
        \item Phân phối (Distributions).
        \item Xu hướng (Trends).
        \item So sánh (Comparisons).
    \end{itemize}
\end{frame}


\begin{frame}{Chọn đúng biểu đồ (Matching Objectives to Types).}
    \begin{itemize}
        \item Thành phần: Pie chart, Stacked bar chart, Waterfall.
        \item Mối quan hệ: Scatterplot, Bubble chart.
    \end{itemize}
\end{frame}


\begin{frame}{Chọn đúng biểu đồ (Tiếp).}
    \begin{itemize}
        \item Phân phối: Histogram.
        \item Xu hướng theo thời gian: Line chart.
        \item So sánh các nhóm: Bar/Column chart.
        ### Phần 3: Đặc điểm của một Data Visualization tốt
    \end{itemize}
\end{frame}


\begin{frame}{Nguyên tắc Nhận thức Thị giác (Visual Perception).}
    \begin{itemize}
        \item Não bộ con người thích sự đơn giản và có trật tự.
        \item Áp dụng các nguyên tắc Gestalt: Continuity, Similarity, Proximity.
    \end{itemize}
\end{frame}


\begin{frame}{Nguyên tắc Tính liên tục (Continuity).}
    \begin{itemize}
        \item Mắt người có xu hướng theo dõi các đường nối hoặc các khối được sắp xếp liên tục.
        \item Ứng dụng: Sắp xếp các cột trong biểu đồ theo thứ tự giảm dần/tăng dần thay vì ngẫu nhiên.
    \end{itemize}
\end{frame}


\begin{frame}{Nguyên tắc Tính tương đồng (Similarity).}
    \begin{itemize}
        \item Các phần tử giống nhau (màu sắc, hình dạng) được hiểu là cùng một nhóm.
        \item Ứng dụng: Dùng chung một màu cho doanh thu của cùng một dòng sản phẩm qua các năm.
    \end{itemize}
\end{frame}


\begin{frame}{Nguyên tắc Khoảng cách (Proximity).}
    \begin{itemize}
        \item Các đối tượng gần nhau được xem là có liên quan mật thiết hơn.
        \item Ứng dụng trong Scatterplot để nhận diện các cụm dữ liệu (Clustering).
    \end{itemize}
\end{frame}


\begin{frame}{Thuộc tính Tiền chú ý (Preattentive Attributes).}
    \begin{itemize}
        \item Là những đặc điểm não bộ nhận diện ngay lập tức trước khi ý thức can thiệp.
        \item Bao gồm: Màu sắc, Kích thước, Hình dạng, Vị trí.
    \end{itemize}
\end{frame}


\begin{frame}{Sử dụng Màu sắc (Color) hiệu quả.}
    \begin{itemize}
        \item Không dùng quá nhiều màu gây rối.
        \item Dùng màu để làm nổi bật (Highlighting) điểm quan trọng.
        \item Chú ý đến người mù màu (Color blindness).
    \end{itemize}
\end{frame}


\begin{frame}{Tránh sự lộn xộn (Avoid Clutter).}
    \begin{itemize}
        \item Xóa bỏ Chart junk (các phần tử trang trí không mang thông tin).
        \item Giảm đường lưới (Gridlines), viền không cần thiết.
        \item Loại bỏ nhãn dữ liệu thừa mứa.
    \end{itemize}
\end{frame}


\begin{frame}{Ứng dụng AI trong việc làm sạch biểu đồ.}
    \begin{itemize}
        \item Prompt AI: "Gợi ý cách tối giản hóa biểu đồ này để tôn lên thông điệp lợi nhuận ròng đang giảm."
    \end{itemize}
\end{frame}


\begin{frame}{Best Practices cho từng loại biểu đồ.}
    \begin{itemize}
        \item Biểu đồ Bar/Column: Trục Y luôn bắt đầu từ 0. Sắp xếp theo một trật tự logic.
    \end{itemize}
\end{frame}


\begin{frame}{Best Practices (Tiếp).}
    \begin{itemize}
        \item Biểu đồ Pie (Tròn): Không dùng quá 5 lát cắt. Không dùng hiệu ứng 3D.
        \item Biểu đồ Line (Đường): Hạn chế số lượng đường (không quá 4-5) để tránh rối mắt.
        ### Phần 4: Nhận diện biểu đồ gây hiểu lầm (Misleading Visualizations)
    \end{itemize}
\end{frame}


\begin{frame}{Tại sao biểu đồ lại lừa dối?}
    \begin{itemize}
        \item Do vô tình thiếu kỹ năng, hoặc cố tình để thao túng người xem (Manipulate).
        \item AI đôi khi cũng có thể sinh ra các biểu đồ sai lệch nếu không được Prompt kỹ.
    \end{itemize}
\end{frame}


\begin{frame}{Lỗi 1: Lược bỏ đường cơ sở (Omitting the Baseline).}
    \begin{itemize}
        \item Biểu đồ cột không bắt đầu từ 0 (Truncated graph).
        \item Phóng đại mức độ chênh lệch giữa các cột.
    \end{itemize}
\end{frame}


\begin{frame}{Lỗi 2: Thao túng trục Y (Manipulating the Y-axis).}
    \begin{itemize}
        \item Thay đổi tỉ lệ (Scale) của trục Y làm cho sự thay đổi trông dốc hơn hoặc bằng phẳng hơn thực tế.
    \end{itemize}
\end{frame}


\begin{frame}{Lỗi 3: Đi ngược lại quy ước thông thường (Going Against Conventions).}
    \begin{itemize}
        \item Ví dụ: Dùng màu Đỏ để chỉ Lợi nhuận tăng, màu Xanh để chỉ Lỗ (trái ngược với quy ước màu sắc kế toán).
        \item Lật ngược trục tọa độ.
    \end{itemize}
\end{frame}


\begin{frame}{Lỗi 4: Chọn lọc dữ liệu (Cherry-picking Data).}
    \begin{itemize}
        \item Chỉ vẽ biểu đồ cho những tháng có doanh thu tăng, bỏ qua các tháng giảm.
        \item Không phản ánh đúng bức tranh toàn cảnh của kỳ kế toán.
    \end{itemize}
\end{frame}


\begin{frame}{Lỗi 5: Dùng sai loại biểu đồ (Using the Wrong Graph).}
    \begin{itemize}
        \item Dùng Pie chart để thể hiện sự thay đổi theo thời gian.
        \item Tổng các phần trong Pie chart không bằng 100\%.
    \end{itemize}
\end{frame}


\begin{frame}{Phản biện AI khi tạo biểu đồ.}
    \begin{itemize}
        \item Luôn kiểm tra xem biểu đồ do AI (Advanced Data Analysis) tạo ra có vi phạm 5 lỗi trên không.
        ### Phần 5: Dashboard và Trực quan hóa tương tác
    \end{itemize}
\end{frame}


\begin{frame}{Thuyết trình trực tiếp với Data Visualizations.}
    \begin{itemize}
        \item Đừng chỉ đọc lại những gì trên màn hình.
        \item Sử dụng kỹ thuật Build-up (Hiển thị từng phần của biểu đồ) để kể chuyện.
    \end{itemize}
\end{frame}


\begin{frame}{Dashboard tương tác (Interactive Visualizations).}
    \begin{itemize}
        \item Người dùng (Quản lý, Nhà đầu tư) có thể tự khám phá dữ liệu.
        \item Tích hợp các bộ lọc (Filters): Theo năm, theo khu vực, theo dòng sản phẩm.
    \end{itemize}
\end{frame}


\begin{frame}{Tích hợp AI vào Dashboard.}
    \begin{itemize}
        \item Sử dụng các công cụ như PowerBI Copilot hoặc Tableau AI để tự động tạo giải thích bằng chữ (Natural Language Generation) bên cạnh biểu đồ.
    \end{itemize}
\end{frame}


\begin{frame}{Ứng dụng trong Báo cáo Quản trị.}
    \begin{itemize}
        \item Thay vì gửi file Excel hàng trăm trang, kế toán gửi một link Dashboard trực tuyến.
        \item Ban giám đốc tự chọn góc nhìn họ cần.
    \end{itemize}
\end{frame}


\begin{frame}{Tổng kết buổi học.}
    \begin{itemize}
        \item Giao tiếp kết quả phân tích là bước quan trọng nhất để tạo ra giá trị.
        \item Nắm vững 3 chữ: Data - Narrative - Visuals.
        \item AI là công cụ tạo biểu đồ, Kế toán là người kiểm soát tính trung thực và thông điệp.
    \end{itemize}
\end{frame}


\begin{frame}{Thực hành - Phần 1.}
    \begin{itemize}
        \item Làm quen với việc phân tích yêu cầu từ sếp và lựa chọn biểu đồ phù hợp cho báo cáo.
    \end{itemize}
\end{frame}


\begin{frame}{Thực hành - Phần 2.}
    \begin{itemize}
        \item Tạo các bảng dữ liệu mẫu, thiết kế Dashboard tĩnh và ứng dụng Data Storytelling để báo cáo.
    \end{itemize}
\end{frame}


\end{document}